% =========================
% chapters/06_system_design_implementation.tex
% =========================
\chapter{System Design and Implementation}
\label{ch:system_design}

\section{Design objectives and constraints}
This project treats the RF control layer as a real-time system. The design objectives are:
(i) deterministic update times relative to RTIO (\texttt{now\_mu}),
(ii) repeatable phase relations across channels and across runs,
(iii) multitone generation for AOM/AOD tasks,
and (iv) an experiment-facing interface that reduces complexity while preserving explicit timing control.

Constraints include: finite tuning word resolution, device-specific update triggers (DAC sync, IO\_UPDATE),
memory/throughput limits for waveform streaming, and analog front-end imperfections (IQ imbalance, filter
responses, amplifier compression).

\section{Software architecture: host orchestration vs kernel scheduling}
Host code is used for configuration, parameter sweeps, data logging, and visualization. Kernel code is used for
time-critical actions, including applying scheduled updates and triggers at deterministic timestamps.

\subsection{Principle: phase control requires a defined reference time}
Phase statements are only meaningful when the reference time $T_0$ is clear. This thesis therefore defines
a ``phase epoch'' and adopts a consistent policy:
\begin{itemize}[leftmargin=*]
  \item Choose a deterministic reference timestamp $T_0$ (e.g., start of experiment segment).
  \item Apply all oscillator/NCO phase settings relative to $T_0$.
  \item Avoid ambiguous ``late updates'' by ensuring positive slack and explicit sync triggers.
\end{itemize}

\section{Multitone output implementation}
\subsection{Tone model and constraints}
% TODO: Describe your chosen gateware mode (classic 5 oscillators, MIQRO 16 tones, or MTDDS 26-tone mode).
% Provide constraints: tone spacing, amplitude scaling, passband limits, clipping/overflow policy.

\subsection{Code skeleton (placeholder)}
\begin{lstlisting}[style=artiq,language=Python,caption={Placeholder: multitone configuration and scheduling.},label={lst:multitone}]
# TODO: Replace with your actual ARTIQ experiment code.
from artiq.experiment import *

class MultiToneDemo(EnvExperiment):
    def build(self):
        self.setattr_device("core")
        self.setattr_device("phaser0")

    @kernel
    def run(self):
        self.core.reset()
        self.core.break_realtime()

        # TODO: init Phaser, set DUC, configure attenuators, and set oscillator tones
        # TODO: schedule updates at deterministic timestamps using now_mu()/at_mu()
\end{lstlisting}

\section{Precision phase control implementation}
This section documents how you enforce phase determinism:
(i) configuration order of operations,
(ii) when and how you clear accumulators,
(iii) whether you use ``tracking/absolute/continuous'' phase modes (as relevant),
and (iv) your calibration of static inter-channel offsets.

% TODO: Insert a timing diagram showing phase-set events and the resulting waveform transitions.
\begin{figure}[t]
  \centering
  \draftfigure[0.95\linewidth]{figures/timing_phase_update.pdf}
  \caption{Placeholder timing diagram: scheduled phase update relative to RTIO timeline cursor and sync trigger.}
  \label{fig:timing_phase_update}
\end{figure}

\section{Experiment-facing interface}
\subsection{Goals}
The interface is designed so experiment scripts express intent (``play this waveform here'') while retaining the
ability to reason about deterministic timing and phase.

\subsection{API outline}
% TODO: Replace with your actual interface and its design rationale.
\begin{lstlisting}[style=artiq,language=Python,caption={Placeholder: experiment-facing interface.},label={lst:interface}]
class RFWaveformInterface:
    """
    TODO: Document: tone specification, waveform scheduling, and metadata logging.
    """
    def __init__(self, phaser_device):
        self.phaser = phaser_device

    def configure(self, config):
        pass

    @kernel
    def apply_at(self, t_mu: int):
        pass
\end{lstlisting}

\section{Validation plan embedded in implementation}
To avoid ``post-hoc validation,'' this project inserts hooks for:
(i) event timestamp logging,
(ii) reference clock identification,
(iii) capture script generation for scope/spectrum analyzer,
and (iv) automated plot generation from captured CSV data.

% TODO: Add your repo layout and data logging conventions.
